\chapter{The World That Doesn't Exist}

\section{What Models Actually Have, Instead of a World}

Chapter 1 described a symptom without asking what, if anything, current systems have in place of the thing that would prevent it. It is worth asking directly, because the honest answer is not nothing. A large language model carries a context window, the recent span of tokens it conditions its next output on, and internally, a set of hidden states or key-value caches summarizing that span in a form suited to fast lookup during generation. A diffusion model carries a noise trajectory and whatever conditioning signal, text, image, or prior frame, it was given for the current generation. A video model carries some window of recent frames it attends to when producing the next one. None of these are nothing. Each is a real, technically sophisticated mechanism, and each does genuinely help a system behave more consistently than a system with no such mechanism at all.

None of them, however, is a world, and this chapter is about the specific way in which they fall short.

\section{The Buffer Is Not the World}

A context window, a hidden state, a cache of recent frames: each of these shares a property worth naming plainly. Each is a buffer. A buffer holds what is currently near at hand, weighted toward recency, and it does not distinguish, in any principled way, between two things that a genuine persistent record would need to keep entirely separate: something that has not happened yet, and something that happened and has since fallen out of the buffer's reach. From inside the buffer's own operation, both cases look identical. Nothing is there. A detail established forty thousand tokens ago and a detail that was simply never mentioned occupy exactly the same status once the conversation has moved on far enough, which is to say no status at all, and a system reasoning only from its buffer cannot tell the two apart because the information needed to tell them apart was never retained anywhere the buffer could still consult it.

This is the wall problem restated in the vocabulary of current architectures rather than the vocabulary of a game engine. The crack that was invented in the first rendering and absent in the second was not maliciously discarded. It fell out of whatever buffer, however defined, the second rendering was conditioned on, and from inside that second rendering's process, there was no way to know a crack had ever been committed to in the first place. A longer buffer pushes this failure further away in time or in scene-distance. It does not remove the failure, because a buffer, by its nature, is finite, and anything a genuine world would need to remember indefinitely will eventually exceed whatever finite span the buffer is willing to hold.

\section{Object Permanence, Borrowed and Misapplied}

Developmental psychology has a name for the achievement current architectures lack, and it is worth borrowing carefully, with an explicit caveat about the borrowing. Human infants, for roughly the first eight to twelve months of life, behave as though an object that has left their field of view has ceased to exist: they do not search for a toy hidden under a cloth in front of them, not because they lack the physical capability to lift the cloth, but because nothing in their current cognitive apparatus represents the toy as continuing to exist somewhere unseen. Object permanence is the developmental achievement of coming to represent objects as persisting independent of being perceived, and its acquisition is one of the most robustly documented milestones in early cognitive development.

The caveat matters more than the analogy. This book is not claiming that language models or diffusion models have minds, that they lack object permanence in the sense an infant lacks it, or that training them longer would constitute development in any psychological sense. The analogy is architectural, not psychological: a system that has no mechanism for representing an unobserved region as continuing to hold specific, determinate content is behaviorally indistinguishable, with respect to that content, from an infant who has not yet acquired object permanence, whatever the underlying explanation for the resemblance turns out to be. Every generation, for a system built entirely from buffers, restarts this developmental stage from zero. There is no accumulation across generations of the kind object permanence, once acquired, provides across a human life. The wall, the crack, the previously established fact about a conversation: each is, from the system's own point of view at the moment of the next generation, an object that has left the field of view and been treated, structurally, as though it no longer exists anywhere at all.

\section{Why This Cannot Be Trained Away}

It is tempting to treat this as a gap that sufficiently targeted training could close: reward consistency, penalize contradiction, and the system should eventually learn to behave as though objects persisted, whether or not anything internal to it actually represents them as persisting. This section argues that training of this kind, however successful within its own terms, addresses the wrong category of problem.

There is a difference between behavior and architecture that this book will insist on throughout. Behavior is a statistical tendency: what a system, on average, across many trials, tends to do. Architecture is a structural fact: what a system is or is not capable of doing, by construction, regardless of how any particular trial happens to go. Training a system to behave more consistently makes contradiction statistically rarer. It does not make contradiction architecturally impossible, and the difference is not merely academic. A system trained to behave consistently will still, on some sufficiently unusual input, sufficiently long context, or sufficiently rare combination of circumstances its training did not adequately cover, produce the very failure the training was meant to suppress, because nothing about the training changed what the system is capable of; it only changed what the system tends to do within the range of situations the training actually sampled. Object permanence, in the architectural sense this chapter needs, cannot be a tendency. It has to be a guarantee, and a guarantee has to come from what the system is built to do, not from what it has merely been encouraged to do often enough.

\section{What Hidden State Was Actually Built For}

It is worth being specific about why hidden states and key-value caches, despite being genuinely sophisticated pieces of engineering, cannot supply this guarantee even in principle, because the reason is not that they are poorly implemented. It is that they were built to solve a different problem. A hidden state in a language model is optimized, through training, to be useful for predicting the next token given everything seen so far. This is an extremely demanding objective, and the representations that emerge from it are correspondingly rich. But usefulness for next-token prediction is not the same property as faithfulness to a specific, checkable, persistent fact. A hidden state can discard information a persistent record would need to keep, provided that information is no longer helpful for predicting what comes next, and a hidden state can blend or approximate information a persistent record would need to keep exact, provided the blend remains statistically useful. Nothing about the training objective ever asked the hidden state to answer the question a genuine memory would need to answer on demand: was this specific fact, this specific detail, already established, and if so, exactly what was established. The hidden state was never built to be queried this way. It was built to be helpful for the next word, and helpfulness for the next word is compatible with silently losing, blending, or reinterpreting exactly the details a persistent world would need to hold fixed.

\section{Looking Ahead}

This chapter has argued that current architectures lack persistence not because their buffers are too short, but because a buffer, however long, is the wrong kind of object to supply what persistence requires, and that this gap cannot be closed by training alone, because training changes behavior, not architecture. It is worth anticipating an objection before the next chapter takes it up directly: some systems already pair a generative model with an explicit external memory, a retrieval index, a vector database, a scratchpad the model can write to and read from. Surely this closes the gap this chapter has described. Chapter 3 examines that possibility directly, and argues that having a memory module is not the same as having a memory that actually functions the way the wall, and everything this book will go on to build, requires.
